\section{Referencial Teórico e Estado da Arte}
\label{sec:referencial}

%A triagem automatizada de vínculos societários ocultos em licitações articula fundamentos da Ciência da Computação (listas de adjacência, conjuntos disjuntos e detecção de motivos em grafos), da Ciência das Redes Complexas e da regulação jurídica das compras públicas. Esta seção consolida o arcabouço normativo, a literatura científica e as escolhas algorítmicas que sustentam o módulo \texttt{01\_\allowbreak analise\_\allowbreak socios}.

\subsection{Paradgima Jurídico-Regulatório e Triagem Indiciária}

As contratações públicas devem assegurar tratamento isonômico entre competidores e seleção da proposta mais vantajosa para a Administração, conforme os arts.~5º, 9º, inciso~I, alínea ``a'', e 11, incisos~I a III, da Lei nº 14.133/2021~\cite{brasil2021licitacoes}. O mesmo diploma tipifica infrações administrativas ligadas à frustração da competitividade e ao conluio (art.~155, incisos~VIII a XII), disciplina a dosimetria sancionatória segundo o dano e a materialidade (art.~156, §~1º), autoriza a desconsideração da personalidade jurídica diante de abuso de direito ou confusão patrimonial entre coligadas (art.~160) e inclui o art.~337-F no Código Penal para tipificar a frustração do caráter competitivo do certame.

No plano concorrencial, a Lei nº 12.529/2011 (art.~36, \textit{caput}, §~3º, inciso~I)~\cite{brasil2011concorrencia} qualifica como infrações à ordem econômica acordos entre concorrentes destinados à fixação de preços ou à divisão de certames. De forma complementar, a Lei Anticorrupção (Lei nº 12.846/2013, art.~5º, inciso~IV)~\cite{brasil2013anticorrupcao} prevê responsabilização objetiva de pessoas jurídicas pela criação de pessoa jurídica interposta ou pela atuação concertada para fraudar licitações, enquanto a Lei Complementar nº 123/2006 (art.~3º, §~4º, incisos~III a V)~\cite{brasil2006lc123} veda a fruição do regime diferenciado de microempresas (ME) e empresas de pequeno porte (EPP) quando há participação societária cruzada acima dos limites legais. Já a Lei de Improbidade Administrativa (Lei nº 8.429/1992, com redação dada pela Lei nº 14.230/2021, arts.~10 e 12)~\cite{brasil1992improbidade} exige demonstração de dolo específico e efetiva lesão ao erário, tornando indispensável o cruzamento entre indícios societários e materialidade executada no orçamento municipal.

Na jurisprudência de controle externo, os Acórdãos nº 2.341/2011 e nº 1.798/2024 do Plenário do TCU~\cite{tcu2011acordao2341,tcu2024acordao1798} firmam diretriz fundamental: a existência isolada de sócios em comum ou parentesco entre licitantes não basta para caracterizar fraude, exigindo convergência com outros elementos objetivos — como co-disputa simultânea, compartilhamento de endereço, telefone, estrutura administrativa e contratação efetiva — para fundamentar apuração em contraditório. Quando os achados advêm de bases públicas identificáveis, incide ainda a LGPD (Lei nº 13.709/2018, arts.~4º, inciso~II, alínea ``b'', 6º, 7º, inciso~IV e §~4º, e 13)~\cite{brasil2018lgpd}, requerendo minimização e pseudonimização de pessoas naturais e jurídicas privadas na divulgação acadêmica.

\subsection{Redes Corporativas, Homofilia e Detecção de Conluio}

Os guias da Organização para a Cooperação e Desenvolvimento Econômico (OCDE)~\cite{oecd2025guidelines} e do Conselho Administrativo de Defesa Econômica (CADE)~\cite{cade2019guia} apontam que um dos principais mecanismos de colusão em compras públicas envolve licitantes ligados por propriedade direta, \textit{holdings}, vínculos familiares ou infraestrutura compartilhada que apresentam lances paralelos ou de cobertura para simular competição~\cite{carbone2024bidrigging}. Dado que tais concorrentes exibem regularidade documental individual em cada certame, indicadores quantitativos apoiados em redes complexas viabilizam a triagem sistemática de grandes acervos públicos~\cite{fazekas2016objective,wachs2020corruption}.

Em estudos empíricos internacionais e nacionais, Wachs e Kertész~\cite{wachs2019cartel} demonstraram que redes de co-participação revelam núcleos densos associados a riscos colusivos em leilões públicos, enquanto Lyra et al.~\cite{lyra2021characterization} mapearam a topologia modular das redes de co-disputa em municípios do Ceará. Avançando sobre conexões societárias, Villamil et al.~\cite{villamil2024collusion} evidenciaram que vínculos indiretos de propriedade e gestão compartilhada amplificam o risco de conluio, e Pompeu e Holanda Filho~\cite{pompeu2025collusion} aplicaram aprendizado em grafos (\textit{GraphSAGE}) ao cruzamento entre licitações brasileiras e QSA. Sob a ótica sociológica, o princípio da homofilia~\cite{mcpherson2001homophily} explica por que atores ligados por laços familiares ou geográficos tendem a formar aglomerados densos, mas também adverte que traços demográficos comuns (como sobrenomes frequentes) produzem associações espúrias se não houver filtragem estatística e cadastral rigorosa. A contribuição deste artigo reside em acoplar o enriquecimento QSA via BrasilAPI a uma cadeia exata de poda determinística ($1$-core, Union-Find e motivos $K_3$), integrada ao filtro anti-falso positivo onomástico e ao saneamento financeiro do SICOM.

\subsection{Grafos Multicamada, Formulação Formal e Escolha Algorítmica}

Diferentemente de projeções monopartidas homogêneas que perdem a distinção semântica entre laços de mercado e laços de propriedade, a literatura de redes multicamada~\cite{kivela2014multilayer,newman2018networks} fornece o formalismo adequado para representar múltiplos tipos de vértices e relações paralelas em um único objeto matemático. Modela-se o ecossistema societário-licitatório como um multigrafo heterogêneo multicamada não direcionado ponderado:
\begin{equation}
  G_{\text{soc}} = (V,\, E,\, \tau,\, \varphi,\, w),
\end{equation}
no qual $\tau: V \to \{\text{Órgão},\, \text{Empresa},\, \text{Sócio Real}\}$ tipifica os vértices, $E$ é o multiconjunto de arestas tipificadas por $\varphi$ nas quatro camadas relacionais (\textit{Participou\_\allowbreak Licitacao}, \textit{Concorreu\_\allowbreak Contra}, \textit{É\_\allowbreak Dono\_\allowbreak De} e \textit{Vinculo\_\allowbreak Suspeito}) — permitindo arestas paralelas de mercado e indício societário entre o mesmo par — e $w: E \to \mathbb{R}_{>0}$ associa o peso semântico (frequência de co-disputa ou escore indiciário).

A seleção dos quatro algoritmos executados em sequência sobre a estrutura própria em listas de adjacência (\texttt{GrafoMulticamada}) justifica-se por eficiência assintótica e interpretabilidade para auditoria:
\begin{enumerate}
  \item \textbf{Subgrafo Induzido por Predicado de Aresta} ($\mathcal{O}(|V|+|E|)$): como os órgãos licitantes funcionam como super-conectores (\textit{hubs} em estrela), a aplicação direta de algoritmos de agrupamento fundiria quase todo o mercado municipal em uma única componente gigante. A indução linear retém apenas as arestas \textit{Vinculo\_\allowbreak Suspeito} ($w \ge 2{,}0$) e as arestas \textit{É\_\allowbreak Dono\_\allowbreak De} incidentes às empresas suspeitas, isolando exatamente o domínio societário de interesse;
  \item \textbf{Decomposição $1$-Core via BFS} ($\mathcal{O}(|V|+|E|)$): fundamentada em Seidman~\cite{seidman1983network}, a poda iterativa em fila de busca em largura remove vértices desconectados ($\text{grau} < 1$) em tempo linear, eliminando empresas e sócios sem vínculos residuais em vez de percorrer subconjuntos redundantes;
  \item \textbf{Componentes via Union-Find} ($\mathcal{O}(|E|\,\alpha(|V|))$): em alternativa a múltiplas varreduras em profundidade, a estrutura de conjuntos disjuntos de Tarjan~\cite{tarjan1975efficiency} — combinando \textit{compressão de caminho} (que achata as árvores ao redirecionar os nós visitados diretamente à raiz durante a consulta) e união por posto (\textit{union by rank}) (que acopla sempre a subárvore de menor altura sob a de maior posto na fusão) — particiona os núcleos conexos ($|C_i| \ge 2$) em tempo praticamente linear, regido pela inversa da função de Ackermann $\alpha(|V|)$;
  \item \textbf{Enumeração de Triângulos $K_3$} ($\mathcal{O}(|E|\,d_{\max})$): em vez de algoritmos genéricos de enumeração de cliques de custo cúbico no pior caso, a interseção direta dos conjuntos pré-indexados de proprietários $\text{Donos}(a) \cap \allowbreak \text{Donos}(b)$ para cada aresta suspeita $(a,b)$ identifica todos os ciclos triangulares fechados (\textit{Sócio} -- \textit{Empresa}$_a$ -- \textit{Empresa}$_b$) proporcionalmente ao grau máximo de propriedade societária $d_{\max}$.
\end{enumerate}

Sob a ótica de métricas de similaridade, descartou-se o coeficiente de Jaccard $J(a,b) = |I_a \cap I_b|/|I_a \cup I_b|$ na ponderação de co-concorrência porque, em pregões com dezenas de itens de consumo, a divisão pela união $|I_a \cup I_b|$ dilui a magnitude absoluta da sobreposição item a item quando uma licitante disputa poucos itens e outra participa do certame integral; por isso, preservam-se as contagens absolutas de certames e itens compartilhados. Igualmente, preteriu-se a ponderação onomástica irrestrita por frequência inversa de documento $\ln(|P|/f(s))$ (IDF logarítmico), pois sobrenomes extremamente raros (ocorrendo em apenas um ou dois pares do município) atingiriam o limiar $\ge 2{,}0$ isoladamente e disparariam falsos positivos no corte estrutural.

Adicionalmente, para diagnosticar a estrutura mesoscópica da macro-rede antes da poda, emprega-se a heurística gulosa de Louvain~\cite{blondel2008fast}, maximizando a modularidade ponderada:
\begin{equation}
  Q = \frac{1}{2m} \sum_{x, y \in V} \left[ A_{xy} - \frac{k_x k_y}{2m} \right] \delta(g_x, g_y),
\end{equation}
em que $A_{xy}$ denota a soma dos pesos entre $x$ e $y$, $k_x = \sum_y A_{xy}$, $m = \frac{1}{2}\sum_{xy} A_{xy}$, $g_x$ é a comunidade de $x$ e $\delta$ é o delta de Kronecker. Para inspeção visual interativa das comunidades e dos núcleos reduzidos no Gephi~\cite{bastian2009gephi}, adota-se o algoritmo contínuo dirigido por forças \textit{ForceAtlas2}~\cite{jacomy2014forceatlas2}, capaz de separar aglomerados densos de vértices periféricos sem distorcer a topologia conectiva.
